\begin{abstract}

Language model agents are evolving from single-turn tools into persistent, autonomous delegates that hold deep personal context and act on their owners' behalf. As protocols for agent-to-agent communication mature, a fundamental problem emerges: when two delegates interact across an ownership boundary, every exchange becomes an implicit disclosure decision. Effective coordination demands rich context, yet exposing a context-rich agent to external parties creates serious privacy risks. We call the set of achievable (security, utility) operating points for a given system configuration the \textit{security--utility frontier}, and its shape is the subject of this thesis.

We make five contributions. First, we \textbf{formalise cross-boundary delegation} as a setting where security enforcement is often pushed into LLM reasoning rather than implemented entirely by deterministic mechanisms, and show why traditional access control models are insufficient on their own. Second, we design and deploy \textbf{SharedOS}, a multi-agent delegation platform with per-relationship context scoping, configurable privacy policies, and tool orchestration. It serves both as a real-world coordination system and as a controlled experimental apparatus. Third, we develop two complementary benchmarks: \textbf{PACT-PAIR} (600 dyadic tasks spanning information and action boundaries, multiple sensitivity categories, relationship-conditioned labels, and prompt-defence gradients) and \textbf{PACT-NET} (997 network tasks over a 25-agent graph, testing third-party provenance, cluster crossings, authority-chain failures, and multi-fact amplification). Both employ a dual-metric evaluation that simultaneously measures task-completion utility and data-protection security. Fourth, we derive \textbf{empirical insights} from these benchmarks: generic policies perform no better than no policy, category-specific policies hit a prompt-level ceiling, multi-turn interaction creates a 3$\times$ incidental leakage gap, and relationship context can shift the cost from leakage to over-refusal. Fifth, motivated by these failures, we evaluate \textbf{architectural interventions}: relationship-specific policy as a behavioural layer, \textit{Mountable Context Cells} as structural data absence, and an \textit{Intelligent Escalation Protocol} that stops ambiguous tool calls before private results enter the model context.

Our findings suggest that prompt-level restriction is necessary but insufficient: robust cross-boundary coordination requires architectural control surfaces that constrain what data is mounted and when tools execute.

\end{abstract}
